\chapter{Station Data}

\section{Role de la station Data}
La station Data fournit la base de verite utilisee par toutes les autres
stations. Sans donnees propres, ni le moteur de signaux ni la validation WFO ne
peuvent produire des conclusions fiables.

\section{Pipeline d'ingestion et de normalisation}
Le pipeline Data couvre :
\begin{enumerate}
\item l'import de donnees (fichiers, refresh provider),
\item la detection/normalisation de format,
\item la fusion incrementale sans doublons,
\item la publication en format canonique OHLCV.
\end{enumerate}

Le Tableau~\ref{tab:data-pipeline} detaille les etapes.

\begin{table}[H]
\centering
\caption{Pipeline de la station Data}
\label{tab:data-pipeline}
\begin{tabular}{p{3cm}p{4cm}p{7.4cm}}
\toprule
\textbf{Etape} & \textbf{Entree} & \textbf{Sortie} \\
\midrule
Ingestion & Excel / provider web & Trames brutes date-prix-volume \\
Normalisation & Colonnes heterogenes & Schema OHLCV unifie \\
Validation & Serie temporelle brute & Serie nettoyee, ordonnee, dedoublonnee \\
Publication & Serie nettoyee & Objet canonique \texttt{market\_data\_store/\{symbol\}/1D} \\
\bottomrule
\end{tabular}
\end{table}

Le Tableau~\ref{tab:data-pipeline} montre que la normalisation est un prealable
technique a toute interpretation de signal.

\section{Fraicheur et qualite des donnees}
La station Data suit deux dimensions : la disponibilite et la fraicheur. Une
metrique operationnelle simple peut etre exprimee par :
\begin{equation}
Q_d = \alpha \cdot C_d + (1-\alpha)\cdot F_d
\label{eq:data-quality}
\end{equation}
ou $C_d$ represente la completude journaliere et $F_d$ la fraicheur relative.
L'equation~\ref{eq:data-quality} sert de score de pilotage interne.

\section{Specificites du marche marocain}
Le contexte marocain impose un traitement explicite :
\begin{itemize}
\item calendrier de jours feries de la place de Casablanca,
\item variations de formats de fichiers historiques,
\item heterogeneite de liquidite entre titres.
\end{itemize}
Ces contraintes sont gerees avant la station Signaux pour eviter
d'interpreter comme signal un artefact de donnees.

\section{Interfaces exposees}
La station Data expose des endpoints de consultation (catalogue, disponibilite,
historique OHLCV) et des endpoints d'action (upload, refresh). Cette separation
limite les effets de bord et facilite le monitoring.

